\begin{table*}[!t]
\centering
\begin{threeparttable}
\caption{Frecuencia de clases de reacción normalizadas (n=200).}
\label{tab:clases}
\footnotesize
\begin{tabular}{lrrrr}
\toprule
Clase & Reacciones & Obtenciones & Total & Proporción \\
\midrule
Oxidación & 34 & 1 & 35 & 17,5\% \\
Hidrogenación & 0 & 26 & 26 & 13,0\% \\
Reducción & 0 & 22 & 22 & 11,0\% \\
Eliminación & 0 & 20 & 20 & 10,0\% \\
Adición & 18 & 0 & 18 & 9,0\% \\
Alquilación & 14 & 2 & 16 & 8,0\% \\
Halogenación & 9 & 0 & 9 & 4,5\% \\
Otras clases & 25 & 29 & 54 & 27,0\% \\
\midrule
Total & 100 & 100 & 200 & 100,0\% \\
\bottomrule
\end{tabular}
\begin{tablenotes}
\footnotesize
\item Las categorías son agrupaciones descriptivas derivadas del texto; no constituyen una caracterización independiente de los mecanismos moleculares.
\end{tablenotes}
\end{threeparttable}
\end{table*}
